% !TeX program = xelatex
% !TeX encoding = UTF-8
% !TeX spellcheck = en_US
% !BIB program = bibtex

%%%%%%%%%%%%%%%%%%%%%%%%%%%%%%%%%%%%%%%%%%%%%%%%%%%%%%%%%%%%%%%%%%%%%%%%%%%%%%%%
%%
%%            JJJJ   K                         K   UUUU         UUUU
%%            JJJJ   KKKK                   KKKK   UUUU         UUUU
%%            JJJJ   KKKKKK               KKKKKK   UUUU         UUUU
%%            JJJJ      KKKKKK         KKKKKK      UUUU         UUUU
%%            JJJJ         KKKKKK   KKKKKK         UUUU         UUUU
%%            JJJJ            KKKKKKKKK            UUUU         UUUU
%%    JJ     JJJJJ               KKK               UUUUU       UUUUU
%%  JJJJJJJJJJJJJ    KKKKKKKKKKKKKKKKKKKKKKKKKKK    UUUUUUUUUUUUUUU
%%    JJJJJJJJJ      KKKKKKKKKKKKKKKKKKKKKKKKKKK      UUUUUUUUUUU
%%
%% This is an example file for using the JKU LaTeX technical report template
%% for your seminar report.
%%
%%
%%%%%%%%%%%%%%%%%%%%%%%%%%%%%%%%%%%%%%%%%%%%%%%%%%%%%%%%%%%%%%%%%%%%%%%%%%%%%%%%

%%%%%%%%%%%%%%%%%%%%%%%%%%%%%%%%%%%%%%%%%%%%%%%%%%%%%%%%%%%%%%%%%%%%%%%%%%%%%%%%
%%
%% Document class: This is a koma-script article.
%%
\documentclass[a4paper,oneside,10pt,ngerman,english]{scrartcl}

%%%%%%%%%%%%%%%%%%%%%%%%%%%%%%%%%%%%%%%%%%%%%%%%%%%%%%%%%%%%%%%%%%%%%%%%%%%%%%%%
%%
%% Treat input files as UTF-8 encoded.
%%
\usepackage[utf8]{inputenc}
%%
%%%%%%%%%%%%%%%%%%%%%%%%%%%%%%%%%%%%%%%%%%%%%%%%%%%%%%%%%%%%%%%%%%%%%%%%%%%%%%%%

%%%%%%%%%%%%%%%%%%%%%%%%%%%%%%%%%%%%%%%%%%%%%%%%%%%%%%%%%%%%%%%%%%%%%%%%%%%%%%%%
%%
%% Use the JKU LaTeX technical report template for this document.
%%
\usepackage[seminarreport,fancyfonts]{jkureport}

%%%%%%%%%%%%%%%%%%%%%%%%%%%%%%%%%%%%%%%%%%%%%%%%%%%%%%%%%%%%%%%%%%%%%%%%%%%%%%%%
%%
%% Load additional packages.
%%

% Using natbib for citet/citep style
\usepackage{natbib}
\setcitestyle{authoryear,round,citesep={;},aysep={,},yysep={;}}

\usepackage{todonotes}
\usepackage{import}
\usepackage{amsfonts}
\usepackage{amsmath}
\usepackage{subfigure}
\usepackage{booktabs}
\usepackage{graphicx}
\usepackage{acronym}

%%
%%%%%%%%%%%%%%%%%%%%%%%%%%%%%%%%%%%%%%%%%%%%%%%%%%%%%%%%%%%%%%%%%%%%%%%%%%%%%%%%


\begin{document}
%%%%%%%%%%%%%%%%%%%%%%%%%%%%%%%%%%%%%%%%%%%%%%%%%%%%%%%%%%%%%%%%%%%%%%%%%%%%%%%%
%%
%% Report information and title page
%%

\title{Seminar in AI - Efficient Self-Verification: Understanding LaSeR for LLM Reasoning}

\subtitle{Seminar Report}

\author{\prefix{} Hannes Brantner \suffix{}\matno{K1614466}}

\submissiondepartment{Institute for Machine Learning}

\abstract{This seminar report critically reviews the work \textit{LaSeR: Reinforcement Learning with Last-Token Self-Rewarding} by \citet{laser}. The study addresses the inefficiency of current self-verification methods in Large Language Models (LLMs) by theoretically deriving that the optimal reasoning reward is equivalent to the model's last-token self-rewarding score. By augmenting standard Reinforcement Learning with Verifiable Rewards (RLVR) with a mean squared error loss on the last token's log-probability, \textit{LaSeR} jointly optimizes reasoning and verification capabilities \citep{laser}. The findings demonstrate that this approach achieves high reasoning accuracy and self-verification performance on mathematical benchmarks with negligible inference cost \citep{laser}. Strengths of the paper include its elegant theoretical foundation and significant efficiency gains. However, limitations exist regarding its generalization to non-mathematical domains and reliance on specific prompt structures \citep{laser}. This report is organized into four sections: a story summary, a paper summary, a critical quality analysis, and an overall conclusion.}

\maketitle
%%
%%%%%%%%%%%%%%%%%%%%%%%%%%%%%%%%%%%%%%%%%%%%%%%%%%%%%%%%%%%%%%%%%%%%%%%%%%%%%%%%

%%%%%%%%%%%%%%%%%%%%%%%%%%%%%%%%%%%%%%%%%%%%%%%%%%%%%%%%%%%%%%%%%%%%%%%%%%%%%%%%
%%
%% Add a table of contents
%%

\cleardoubleoddpage
\tableofcontents

%%
%%%%%%%%%%%%%%%%%%%%%%%%%%%%%%%%%%%%%%%%%%%%%%%%%%%%%%%%%%%%%%%%%%%%%%%%%%%%%%%%

%%%%%%%%%%%%%%%%%%%%%%%%%%%%%%%%%%%%%%%%%%%%%%%%%%%%%%%%%%%%%%%%%%%%%%%%%%%%%%%%
%%
%% Story Summary
%%

\cleardoubleoddpage

\section{Story Summary}
\label{sec:storysummary}

\subsection*{1. What is the central question?}
The paper addresses a critical efficiency bottleneck in Large Language Models (LLMs): How can reasoning and self-verification capabilities be unified within a single model to provide reliable signals at test time, without incurring the high computational costs associated with external verifiers or sequential self-correction? \citep{laser}.

\subsection*{2. Why is this question important?}
Reinforcement Learning with Verifiable Rewards (RLVR) typically relies on ground-truth answers, which are unavailable during test-time inference \citep{laser}. While verifiable answers are accessible in deterministic domains like mathematics, they are elusive in subjective domains (e.g., creative writing or summarization) where formal correctness is difficult to assess. Although methods such as external verifiers or self-correction prompts exist, they are inefficient: external verifiers require training separate models, while sequential self-verification doubles the inference cost \citep{laser}. Solving this problem enables efficient \textit{test-time scaling}, allowing models to self-select correct answers with minimal computational overhead \citep{laser}.

\subsection*{3. What evidence/data (variables) are needed to answer this question?}
To address this question, the authors must measure reasoning accuracy (e.g., Pass@1), self-verification performance (F1 score), and computational cost (inference tokens) \citep{laser}. Furthermore, they must validate the theoretical correlation between the \textit{last-token self-rewarding score} (derived from a log-probability ratio) and the true verifier reward \citep{laser}, demonstrating that this ratio can be computed via a simplified, efficient form.

\subsection*{4. What methods are used to get this evidence/data?}
\citet{laser} propose \textit{LaSeR} (Reinforcement Learning with Last-Token Self-Rewarding), an algorithm that augments the standard RLVR objective with a Mean Squared Error (MSE) loss. This loss aligns the model's last-token log-probability for a specific token (e.g., \texttt{<vision\_start>}) with the solution's true correctness \citep{laser}. The authors evaluate this method against baselines such as GRPO and external Reward Models (PRMs/ORMs) using architectures based on LLaMA-3.2 and Qwen-2.5 \citep{laser}.

\subsection*{5. What analyses must be applied for the data to answer the central question?}
To validate the methodology, the authors must empirically demonstrate that the simplified theoretical assumption, specifically that the partition function $Z(x,y) \approx 1$, holds true \citep{laser}. Additionally, they must analyze the correlation between the learned self-rewarding scores and ground-truth correctness across various mathematical benchmarks, such as MATH500 and AIME \citep{laser}.

\subsection*{6. What evidence/data (values for the variables) were obtained?}
\textit{LaSeR} achieved reasoning accuracy comparable to, or exceeding, standard GRPO (e.g., 80.2\% vs. 79.9\% on MATH500 for Qwen2.5-7B) \citep{laser}. Crucially, the method demonstrated high self-verification F1 scores (approximately 80\%) \citep{laser}, matching the performance of a 72B external reward model while requiring only one additional token generation per solution \citep{laser}. Notably, these self-assessment capabilities remained robust even on difficult datasets like AIME-24, where reasoning accuracy was significantly lower (around 15\%) \citep{laser}. This represents a substantial advancement in the model's ability to distinguish correct from incorrect outputs at test time, effectively mitigating issues related to uncalibrated confidence.

\subsection*{7. What were the results of the analyses?}
The analysis confirmed that the \textit{last-token self-rewarding score} effectively captures solution correctness, enabling the model to act as its own verifier during both training and testing \citep{laser}. The simplification of the reference model probability to a constant was empirically validated, which further reduced computational overhead \citep{laser}. However, the method's generalization to \textit{inexact} domains (such as general reasoning tasks in MMLU-Pro) was observed to be less effective compared to the mathematical domain \citep{laser}.

\subsection*{8. How did the analyses answer the central question?}
The results demonstrate that reasoning and verification can be jointly optimized without significant overhead \citep{laser}. By encoding the verification signal directly into the last token's probability distribution, \textit{LaSeR} eliminates the need for a secondary inference pass, successfully unifying the roles of generator and verifier \citep{laser}.

\subsection*{9. What does this answer tell us about the broader field?}
These findings imply that LLMs possess latent self-evaluation capabilities that can be unlocked via specific RL objectives, without the need for complex architectural changes or massive external reward models \citep{laser}. This suggests a viable path toward efficient \textit{System 2} reasoning, where models inherently recognize the validity of their outputs \citep{laser}. While the concept of models quantifying uncertainty is established in Explainable AI (XAI), this work advances the field by integrating it directly into the generation process for efficient reinforcement learning.

\subsection*{10. Did the paper answer the question satisfactorily? Why (not)?}
Yes, the paper satisfactorily answers the central question by providing a rigorous theoretical derivation supported by extensive empirical results on standard benchmarks \citep{laser}. The method's efficiency, requiring only a single additional token, directly resolves the core challenge of computational cost in self-verification \citep{laser}.

%%%%%%%%%%%%%%%%%%%%%%%%%%%%%%%%%%%%%%%%%%%%%%%%%%%%%%%%%%%%%%%%%%%%%%%%%%%%%%%%
%%
%% Paper Summary
%%

\section{Paper Summary}
\label{sec:papersummary}

This chapter summarizes the paper \textit{LaSeR: Reinforcement Learning with Last-Token Self-Rewarding} by \citet{laser}. The summary is structured to reflect the logical flow of a machine learning research paper, covering the introduction, methodology, experimental setup, and results.

\subsection{Introduction}
The landscape of Large Language Models (LLMs) is undergoing a paradigm shift towards \textit{System 2} reasoning capabilities, where models are expected to engage in deliberate thought processes rather than simple pattern matching. Reinforcement Learning with Verifiable Rewards (RLVR) has emerged as a cornerstone for this advancement, exemplified by systems like OpenAI o1 and DeepSeek-R1 \citep{r1}. RLVR encourages models to produce deliberate reasoning chains by rewarding final outcomes based on strict ground-truth consistency. However, a significant operational dilemma exists: while RLVR is effective during training where labels exist, it fails to provide verification signals during test-time inference when ground-truth answers are unavailable \citep{laser}.

To address this \textit{test-time blind spot}, prior research has relied on two primary strategies: training external verifiers (Reward Models) or utilizing sequential self-verification (e.g., asking the model to \textit{double-check} its answer). \citet{laser} argue that both approaches suffer from severe inefficiency. External verifiers impose a high computational burden by requiring the maintenance and query of separate, often parameter-heavy models. Sequential self-verification, while unifying the model, effectively doubles the inference cost and latency by necessitating a second full generation pass to critique the first \citep{laser}. In response, the authors introduce \textit{LaSeR} (Reinforcement Learning with Last-Token Self-Rewarding), a novel algorithm designed to unify reasoning and verification within a single policy model at nearly zero additional cost \citep{laser}.

\subsection{Methods}
\textbf{Theoretical Framework.} The core contribution of the paper is the theoretical derivation linking the RL objective of verification to a closed-form solution. Standard Reinforcement Learning from Human Feedback (RLHF) aims to maximize the expected reward while maintaining proximity to a reference policy $\pi_{\text{ref}}$. The objective function is formally defined as:
\begin{equation}
	\max_{\pi_\theta} \mathbb{E}_{x \sim \mathcal{D}, y \sim \pi_\theta(y|x)} \left[ r(x, y) - \beta D_{\text{KL}}(\pi_\theta(y|x) || \pi_{\text{ref}}(y|x)) \right]
\end{equation}
where $r(x, y)$ is the reward for prompt $x$ and response $y$, and $\beta$ controls the KL-divergence penalty. As derived in \citet{directpreferenceopt}, the optimal policy $\pi_r$ maximizing this objective takes the closed form:
\begin{equation}
	\pi_r(y | x) = \frac{1}{Z(x)} \pi_{\text{ref}}(y | x) \exp\left(\frac{1}{\beta} r(x, y)\right)
\end{equation}
where $Z(x)$ is the partition function. \citet{laser} apply this relation to the prediction of a confidence token $z_c$ (e.g., a special unused token like \texttt{<vision\_start>}) conditioned on the context $x$ and the generated reasoning $y$. Substituting the model's policy $\pi_\theta(z_c | x, y)$ for $\pi_r$, the reward derivation becomes:
\begin{equation}
	r(x, y) = \beta_v \log \frac{\pi_\theta(z_c | x, y)}{\pi_{\text{ref}}(z_c | x, y)} + \beta_v \log Z(x, y)
\end{equation}
\citet{laser} assume that because $z_c$ is an unused special token, its probability mass is negligible, and thus the partition function $Z(x, y)$ approximates to 1, allowing the term $\beta_v \log Z(x, y)$ to be discarded. Furthermore, empirical observation suggests that the reference policy probability $\pi_{\text{ref}}(z_c | x, y)$ for such tokens is effectively constant. Applying these simplifications yields the final \textit{LaSeR} expected reward formulation:
\begin{equation}
	r_{\text{LaSeR}}(x, y) = \beta_v \log \pi_\theta(z_c | x, y) - \beta_v c_{\text{ref}}
\end{equation}
where $\beta_v$ is a scaling coefficient and $c_{\text{ref}} \approx \log \pi_{\text{ref}}(z_c | x, y)$ is the constant baseline. Since probabilities lie in $(0, 1]$, the logarithmic term implies that this reward signal is theoretically unbounded in the negative direction \citep{laser}. This insight allows the authors to bypass complex implicit reward calculations, converting a global optimization problem into a local regression task.

\textbf{LaSeR Algorithm and Optimization.} Based on this theoretical insight, \citet{laser} propose a lightweight augmentation to the standard RLVR process. \textit{LaSeR} adds a Mean Squared Error (MSE) loss to the primary RL objective. This auxiliary loss forces the \textit{last-token self-rewarding score} (the scaled log-probability difference) to regress towards the actual verifier reward (0 or 1) derived from the ground truth. This effectively \textit{distills} the external verifier's judgment into the probability distribution of a single token at the end of the generation \citep{laser}.

To ensure stable training, \citet{laser} introduce several crucial techniques. First, they employ a \textbf{class-level loss re-weighting} strategy to handle the imbalance between correct and incorrect solutions, preventing the model from biasing its verification towards the majority class. Second, they utilize \textbf{Self-Rewarding-based Advantages (SWA)}, where the learned self-rewarding score is fed back into the RL advantage estimation. This creates a dense reward signal, smoothing out the sparse binary feedback of the hard verifier \citep{laser}. To prevent fitting to noise, this intrinsic signal is gated based on the \textit{confidence standard deviation}; it contributes to the advantage only when the model demonstrates sufficient discrimination between samples in a group. Additionally, the protocol includes a \textit{warmup phase} where intrinsic rewards are initially suppressed to prevent policy collapse due to poor initial calibration \citep{laser}.

\textbf{Efficiency via Simplification.} A significant methodological innovation is the simplification of the reference model term. \citet{laser} empirically observe that the log-probability of the special token under the reference model, $\log \pi_{ref}(z_c|x,y)$, is effectively a constant $c_{ref}$ across different inputs due to the token's irrelevance in standard pre-training. By treating this as a pre-calculated constant, \textit{LaSeR} eliminates the need for a forward pass through the reference model during the self-rewarding calculation, reducing the computational overhead to merely one additional token inference \citep{laser}.

\subsection{Experiments}
\textbf{Setup and Baselines.} The framework was evaluated using LLaMA-3.2 (3B) and Qwen-2.5 (7B) architectures trained on the DeepMath-103K dataset \citep{deepmath103k}. \citet{laser} benchmarked \textit{LaSeR} against three distinct baselines: standard pre-trained base models, models trained with standard GRPO (Group Relative Policy Optimization) without self-rewarding, and varying sizes of external Reward Models (ORMs and PRMs), including a 72B parameter model \citep{laser}.

\textbf{Datasets and Metrics.} Evaluation was conducted across five rigorous mathematical reasoning benchmarks: MATH500, AMC23, AIME24, AIME25, and OlympiadBench. The primary metrics assessed were Pass@1 accuracy for reasoning capabilities and the F1 score for self-verification performance, ensuring the model could accurately identify both correct and incorrect solutions \citep{laser}.

\subsection{Results}
\textbf{Reasoning Performance.} Contrary to the concern that multi-tasking might degrade performance (the \textit{alignment tax}), \textit{LaSeR} consistently matched or outperformed the GRPO baseline. For instance, on the Qwen2.5-7B architecture, \textit{LaSeR} achieved an average accuracy of 53.5\% across benchmarks compared to GRPO's 49.2\% \citep{laser}. \citet{laser} attribute this to the improved confidence calibration and the integration of self-rewarding advantages (SWA) during training, which provides a smoother gradient signal than sparse binary rewards.

\textbf{Self-Verification and Scaling.} The most striking result involves \textit{inference-time scaling}. \textit{LaSeR} models achieved self-verification F1 scores of approximately 80\%, performing on par with the Qwen2.5-Math-RM-72B external verifier \citep{laser}. This high-fidelity self-score allows for weighted majority voting, where the model's own confidence scores are used to rank and select the final answer from multiple sampled paths. As shown in Figure 2 of the paper, this weighted voting significantly outperforms standard majority voting, unlocking higher performance with the same inference budget \citep{laser}.

%%%%%%%%%%%%%%%%%%%%%%%%%%%%%%%%%%%%%%%%%%%%%%%%%%%%%%%%%%%%%%%%%%%%%%%%%%%%%%%%
%%
%% Paper Analysis
%%

\section{Paper Analysis}
\label{sec:paperanalysis}

This section provides a critical analysis of the work by \citet{laser}, evaluating the quality of its content, the robustness of its methodology, and the implications of its findings for the broader AI research community.

\subsection{Content Analysis}
\textbf{Theoretical Novelty and the \textit{Token Hijacking} Mechanism.} The research motivation behind \textit{LaSeR} is well-founded in the current bottleneck of LLM deployment: the high cost of \textit{System 2} verification. The theoretical contribution by \citet{laser}, which equates the verifier reward to the last-token log-probability, constitutes an elegant reduction. It transforms what is typically a complex auxiliary task (training a separate reward model) into a simple regression task on the policy model's logits. A particularly clever aspect is the utilization of unused vocabulary tokens (like \texttt{<vision\_start>}). This \textit{token hijacking} effectively repurposes the massive residual capacity of the LLM's vocabulary space to carry semantic metadata (correctness) without interfering with the natural language generation generation \citep{laser}. The derivation that the partition function $Z(x,y)$ can be ignored is insightful, but it relies heavily on the assumption that the reference model maintains near-zero probability for these tokens. While empirically validated here, this assumption represents a potential point of fragility if base models are updated or if token distributions shift significantly in future architectures \citep{laser}.

\textbf{Practical and Structural Limitations.} The \textit{LaSeR} mechanism relies on generating an additional token \textit{after} the \texttt{EOS} (End of Sequence) token. This is a non-standard inference mechanism. Standard inference engines are typically engineered to terminate generation immediately upon encountering \texttt{EOS}. Consequently, adopting \textit{LaSeR} requires modifying low-level inference loops, which may hinder \textit{plug-and-play} adoption in existing production pipelines or inference libraries and frameworks. Furthermore, by compressing all prior reasoning uncertainty into a single final step, the method may be susceptible to length bias in complex chains. In long reasoning trajectories, the probability distribution of the final token may not accurately reflect logic errors that occurred significantly earlier in the chain. The mechanism remains an aggregation heuristic that might fail in highly complex, multi-stage reasoning tasks where errors are subtle and distributed.

\textbf{The Generator-Verifier Gap and Calibration.} A critical paradox in self-verification research is the \textit{Generator-Verifier Gap}: if a model is capable of verifying that an answer is incorrect, why did it generate that answer in the first place? \textit{LaSeR} addresses this implicitly through calibration. By forcing the model to predict a correctness score \textit{immediately} after generation, \textit{LaSeR} is essentially training the model to recognize its own uncertainty \citep{laser}. The superior performance of \textit{LaSeR} over GRPO suggests that this explicit calibration signal helps the model internalize failure modes during training. However, the method is fundamentally bound by the capabilities of the ground-truth verifier used during training. In mathematical domains, this verifier is absolute. In general reasoning, as shown by the drop in performance on MMLU-Pro, the noise in the training verifier propagates directly into the self-rewarding score \citep{laser}. This indicates that \textit{LaSeR} is an excellent amplifier of verifiable truth but cannot synthesize verification capabilities ex nihilo.

\textbf{RL Dynamics: Dense vs. Sparse Rewards.} The introduction of Self-Rewarding-based Advantages (SWA) is a subtle but profound contribution to the Reinforcement Learning dynamics. Standard RLVR typically utilizes sparse binary rewards (0 for incorrect, 1 for correct). Sparse rewards are notoriously difficult for gradient descent optimization due to the credit assignment problem. By mixing the model's continuous self-score into the advantage function, \textit{LaSeR} effectively densifies the reward landscape \citep{laser}. This continuous signal allows the model to differentiate between \textit{confident errors} and \textit{near misses}, likely explaining the accelerated reasoning performance even without test-time verification. This aligns with broader findings in RL that shaping rewards leads to faster convergence, but \textit{LaSeR} achieves this shaping endogenously rather than requiring manual reward engineering.


\textbf{Prompt Sensitivity.} While efficient, the method relies on specific prompt structures to trigger the verification state. The experiments show high sensitivity to prompts, a common trait in LLMs \citep{laser}. This is because the external verifier needs to be presented with the answer in a consistent format in order to extract the answer deterministically for correctness evaluation. This answer format is given to the model with the system prompt.

\subsection{Scientific Writing \& Layout Analysis}
\textbf{Structure and Clarity.} The paper demonstrates high-quality scientific writing. The separation between the theoretical derivation (Section 3) and the practical algorithm (Algorithm 1) aids readability, allowing readers to understand the \textit{why} before the \textit{how} \citep{laser}. The authors are intellectually honest about their results; specifically, the comparison of self-rewarding score distributions (Figure 3) provides clear visual evidence of the separability in Math tasks versus the significant overlap in General Reasoning tasks. This visual honesty regarding the method's limitations strengthens the credibility of the positive results \citep{laser}.

\textbf{Reproducibility.} The paper excels in reproducibility. The appendix provides exhaustive details on hyperparameters (learning rates, batch sizes, $\beta$ coefficients) and, crucially, the specific prompt templates used \citep{laser}. Given that RL performance is often hypersensitive to prompting and initialization, sometimes referred to as the \textit{seed lottery}, this level of detail is commendable and necessary for the community to build upon this work \citep{laser}.

\section{Conclusion}
\label{sec:conclusion}

In this work, \citet{laser} propose \textit{LaSeR}, a framework that successfully democratizes verification capabilities for Large Language Models. By theoretically proving that the optimal reasoning reward can be compressed into the log-probability of a single token, the authors achieve a unification of \textit{actor} and \textit{critic} that was previously computationally prohibitive \citep{laser}.

The implications of this research extend beyond simple efficiency. \textit{LaSeR} demonstrates that expensive, external 72B-parameter reward models can be effectively approximated by a 7B policy model's own latent confidence scores, provided the training objective is correctly aligned \citep{laser}. This enables a new paradigm of \textit{inference-time scaling}, where models can generate multiple solutions and self-select the best one with minimal latency overhead. The shift from computing implicit rewards over the entire sequence to a single-token check represents a move towards \textit{Compute-Optimal} inference, shifting the heavy lifting from runtime verification to training-time calibration \citep{laser}.

Future work will likely focus on closing the generalization gap observed in non-mathematical domains. Future directions may also include removing the dependence on the single additional forward pass after the \texttt{EOS} token to streamline integration. Furthermore, there is potential to extend the current \textit{last token} verification into a finer-grained \textit{step by step} verification, where a self-reward signal is generated at each token generation step. Overall, \textit{LaSeR}represents a decisive step towards autonomous, self-correcting AI systems that are both capably \textit{System 2} and computationally viable for widespread deployment \citep{laser}. The additional self-rewarding score further enables LLMs to better quantify their uncertainty, contributing to the broader goals of Explainable AI (XAI) by providing interpretable confidence metrics alongside generated outputs.

%%
%%%%%%%%%%%%%%%%%%%%%%%%%%%%%%%%%%%%%%%%%%%%%%%%%%%%%%%%%%%%%%%%%%%%%%%%%%%%%%%%

%%%%%%%%%%%%%%%%%%%%%%%%%%%%%%%%%%%%%%%%%%%%%%%%%%%%%%%%%%%%%%%%%%%%%%%%%%%%%%%%
%%
%% Print the bibliography
%%
\cleardoubleoddpage
\bibliographystyle{plainnat}
\bibliography{references}
%%
%%%%%%%%%%%%%%%%%%%%%%%%%%%%%%%%%%%%%%%%%%%%%%%%%%%%%%%%%%%%%%%%%%%%%%%%%%%%%%%%

\end{document}